% !TeX root = ../机器学习笔记.tex
% 由主文档通过 \input 载入；请编译主文档。
\chapter*{绪论：课程总纲与学习路线}
\addcontentsline{toc}{chapter}{绪论：课程总纲与学习路线}

\section*{研究对象与全书定位}

给定若干房屋的面积与成交价格，如何预测另一套房屋的价格；给定一批邮件及其类别，如何判断新邮件是否为垃圾邮件；只有用户行为记录而没有类别时，又能发现什么结构？这些问题共同要求从有限观测中构造可用于新情形的规则。机器学习的基本问题因此包括如何表示数据、如何选择模型、如何利用数据确定模型，以及如何判断所得规律是否可靠。本书以“任务与数据---模型与目标---优化与计算---评价与泛化”为贯穿路线，强调算法背后的假设、推导和适用条件。

本书面向具备初步微积分、线性代数和概率知识的本科生，先建立经典机器学习的完整主干，再介绍神经网络及相关进阶方向。数学工具在使用前集中回顾；抽象理论以能够解释实际建模选择为边界，涉及更深结果时区分完整证明、证明思路和选读内容。神经网络和强化学习在这里承担课程衔接作用，深入专题可继续在相应学科笔记中展开。

本书采用紧凑的本科教材体例：基础概念给出定义，核心结论交代条件并提供关键证明，算法明确目标、更新和停止规则，例题用于贯通推导与计算。进阶专题侧重机制与适用范围，较深的引用结论不冒充已经证明的结果。章末实验是供读者实施的学习任务，书中不将尚未执行的实验写成已有结果。

\section*{课程总纲}

全书分五篇、二十五章。前四篇形成本科主干；第五篇中专题可按兴趣选读，综合实验则应与主干学习同步开展。

\begin{description}
  \item[第一篇：学习问题与数学基础]
  从任务、数据和反馈出发，建立理解模型、训练和评价所需的共同语言。
  \begin{enumerate}
    \item 第\ref{chap:ml-tasks}章：机器学习的任务与基本概念。
    \item 第\ref{chap:ml-linear-algebra}章：线性代数与微积分工具。
    \item 第\ref{chap:ml-probability}章：概率、统计与信息论基础。
    \item 第\ref{chap:ml-optimization}章：机器学习中的优化方法。
    \item 第\ref{chap:ml-evaluation}章：数据处理、模型评估与选择。
    \item 第\ref{chap:ml-generalization}章：泛化、过拟合与正则化。
  \end{enumerate}
  \item[第二篇：监督学习的基本方法]
  按线性模型、概率模型、局部方法、树模型与核方法的顺序，比较不同的归纳假设。
  \begin{enumerate}
    \item 第\ref{chap:ml-regression}章：线性回归与正则化回归。
    \item 第\ref{chap:ml-logistic}章：感知机、逻辑回归与最大熵。
    \item 第\ref{chap:ml-generative}章：贝叶斯决策与概率生成分类。
    \item 第\ref{chap:ml-neighbors}章：近邻方法与非参数估计。
    \item 第\ref{chap:ml-trees}章：决策树。
    \item 第\ref{chap:ml-ensembles}章：集成学习。
    \item 第\ref{chap:ml-kernels}章：核方法与支持向量机。
  \end{enumerate}
  \item[第三篇：无监督学习与隐变量模型]
  从几何分组、低维近似和概率建模三个角度认识数据结构。
  \begin{enumerate}
    \item 第\ref{chap:ml-clustering}章：聚类分析。
    \item 第\ref{chap:ml-dimension}章：降维、主成分分析与矩阵分解。
    \item 第\ref{chap:ml-em}章：隐变量模型与期望最大化算法。
  \end{enumerate}
  \item[第四篇：神经网络与表示学习]
  将前面关于损失、优化和泛化的认识推广到可学习的非线性特征。
  \begin{enumerate}
    \item 第\ref{chap:ml-networks}章：前馈神经网络与反向传播。
    \item 第\ref{chap:ml-representation}章：深度模型与表示学习初步。
  \end{enumerate}
  \item[第五篇：进阶专题与综合实践]
  围绕结构化数据、不完全监督、序贯决策及可靠实验组织选读内容。
  \begin{enumerate}
    \item 第\ref{chap:ml-graphical}章：概率图模型与序列学习。
    \item 第\ref{chap:ml-bayesian}章：贝叶斯学习与近似推断。
    \item 第\ref{chap:ml-text}章：文本表示与主题模型。
    \item 第\ref{chap:ml-graphs}章：图数据、推荐与排序。
    \item 第\ref{chap:ml-transfer}章：半监督、主动学习与迁移学习。
    \item 第\ref{chap:ml-rl}章：强化学习导论。
    \item 第\ref{chap:ml-practice}章：综合实验与可靠机器学习实践。
  \end{enumerate}
\end{description}

附录\ref{app:ml-notation}整理记号与常用计算，附录\ref{app:ml-methods}提供方法比较与学习自检，附录\ref{app:ml-materials}给出本书与现有分章材料的对应关系。章节按知识依赖组织，已有材料的编号与本书章号分别使用。

\section*{先修知识与学习顺序}

初学者先读第\ref{chap:ml-tasks}章，并在第\ref{chap:ml-linear-algebra}至第\ref{chap:ml-optimization}章补足所需工具，再学习第\ref{chap:ml-evaluation}章；之后可先完成第\ref{chap:ml-regression}章的一个回归例子，再返回第\ref{chap:ml-generalization}章理解泛化。第\ref{chap:ml-neighbors}章和第\ref{chap:ml-trees}章所需推导较少，可作为较早的算法实践。每学习一种模型，都回到同一数据划分和评价规则下与简单基线比较。

支持向量机依赖凸优化与对偶；主成分分析依赖谱分解；期望最大化依赖条件概率、似然和 Jensen 不等式；反向传播依赖矩阵求导与链式法则。概率图模型之后可继续学习序列模型和近似推断；主题模型依赖矩阵分解及隐变量建模；推荐依赖矩阵分解，图排序依赖有限状态马尔可夫链。强化学习在任务上引入行动对数据分布的影响，不应仅视为另一种静态分类。

\section*{全书记号与论证约定}

随机输入和输出记为 $X$、$Y$，观测值记为 $x$、$y$。监督训练集写为 $\mathcal D=\{(x_i,y_i)\}_{i=1}^{n}$；输入通常为 $x_i\in\mathbb R^d$，其中 $n$ 为样本数，$d$ 为特征数。样本设计矩阵用粗体 $\mathbf X\in\mathbb R^{n\times d}$ 表示，其第 $i$ 行是 $x_i^\top$，与随机向量 $X$ 区分。参数记为 $\theta$，预测函数记为 $f_\theta$，预测值记为 $\widehat y$，损失统一按 $\ell(\widehat y,y)$ 的次序书写。需要截距时明确单独表示或并入扩展特征，不混用两种约定。

一般分类标签使用 $\{1,\ldots,C\}$，其中 $C$ 为类别数；二分类逻辑回归使用 $\{0,1\}$，感知机与支持向量机使用 $\{-1,+1\}$，转换时明确说明。聚类的簇数记为 $K$，近邻数记为小写 $k$。正则化强度记为 $\lambda\geq0$，学习率记为 $\eta>0$；这些通常是由训练流程外部设定或选择的超参数，不与直接拟合的参数混淆。

未另行说明时，基础统计分析假设训练样本独立同分布，测试样本独立于训练过程并来自同一分布；时间序列、用户重复观测和分布变化另行处理。证明必须说明随机性来自抽样、噪声还是算法；“目标值不增”“迭代收敛”“到达全局最优”和“泛化良好”是不同性质，分别论证。

\section*{阅读与练习}

各章由问题引出概念，依次讨论模型、目标、推导、算法及评价，章末给出手算例题与练习。理论结论先规定假设，实验先规定数据划分和评价指标；手算用于检查机制，实验用于观察有限数据下的行为，二者均不能替代一般定理的证明。算法使用与软件无关的数学描述；读者采用具体库实现时，应记录版本并核对接口、默认参数和数值约定。

